\documentclass{report}
\usepackage{amsfonts, amsmath}
\newcommand\R{{\mathbb R}}
\newcommand\Q{{\mathbb Q}}
\newcommand\N{{\mathbb N}}
\pagestyle{empty}
\begin{document}
\large
\centerline{\Huge Analisi Matematica II}
\centerline{\large Corso di Laurea di Informatica}
\renewcommand{\thefootnote}{ } 
\footnote{\sl Avete 2.30 ore di tempo. Ogni esercizio
vale nove punti. La sufficienza si ottiene con un punteggio $\geq$
18. Solo le risposte chiaramente giustificate saranno prese in
considerazione.}
\renewcommand{\thefootnote}{\arabic{footnote}} 
\vskip.1cm
\centerline{\Large Appello 8-09-2003}
\vskip1cm
\begin{enumerate}
\item Si scriva la serie di Taylor, in zero, della funzione
\[
f(x)=\sqrt{1-x}
\]
e se ne studi la convergenza.
\item Si mostri che
\[
\sqrt{1-x^2}\leq 1-\frac 12x^2-\frac 18 x^4\quad \forall x\in [-1,1].
\]
Si usi tale fatto per dimostrare che $\pi\leq 3.24$. Sapreste fare di meglio?
\item Per ogni $\gamma>2$ si risolva il problema di Cauchy
\[
\begin{split}
&\ddot x=-\gamma\dot x-x\\
&x(0)=1;\;\; \dot x(0)=0.
\end{split}
\]
Detta $x_\gamma(t)$ la soluzione di tale problema si calcoli
\[
\lim_{\gamma\to\infty}x_\gamma(t).
\]
\item Si mostri che
\[
f(x)=\frac 12\int_{x-1}^{x+1}|y|dy
\]
\`e ovunque derivabile e se ne calcoli la derivata in zero.
\end{enumerate}
\end{document}
